\begin{activity}{Inversion of a Buried Sphere}

\section*{Purpose}

The purpose of this activity is to develop an intuitive understanding
of inversion using the gravity anomaly of a buried sphere.
By comparing grid-search and gradient-based approaches, you will
explore how inverse methods search model space, how model
parameterization influences the solution, and why non-uniqueness and
noise make geophysical interpretation challenging. 

\section*{Learning Outcomes}

By the end of this activity, you should be able to:
\begin{itemize}
    \item Explain the difference between forward modelling, grid search,
          and gradient-based inversion.
    \item Describe the role of the Jacobian matrix in determining model updates.
    \item Identify sources of non-uniqueness in geophysical inverse problems.
    \item Explain how model parameterization can improve or degrade
          inversion performance.
    \item Describe the influence of noise on inversion results and
          uncertainty.
\end{itemize}

\section*{Gravity of a Buried Sphere}

We will return to how we derive the gravity due to a buried sphere later.  For now, we will model the vertical gravity anomaly at the surface produced by a buried sphere of depth $z$ (to center), radius $a$, and density contrast $\Delta\rho$. For a station at horizontal radius $r=\sqrt{x^2+y^2}$,
\begin{equation}
  g(r|z_c,a,\Delta\rho) = \underbrace{\frac{4}{3} \pi a^3 \Delta\rho}_{\text{mass}} G \frac{z}{\left( r^2 + z_c^2 \right)^{3/2}}\, ,
  \label{eq:forward}
\end{equation}
with $G$ the gravitational constant. 

\section*{Inverse Model}

The gravity equation due to a buried sphere is non-linear, requiring a non-linear iterative solver. The code we will use employs the Gauss--Newton inverse method, which can be enhanced with Levenberg--Marquardt regularization.

We seek $\vb{m} = [z_c, a, \Delta\rho]\T$ that minimizes the objective function,
\[
  \Phi(m) = \underbrace{\left\| W_d\,(\vb{g}_{\mathrm{obs}} - g(\vb{m})) \right\|_2^2}_{\text{model misfit}} + \underbrace{\lambda^2\,\left\| (\vb{m} - \vb{m}_{0}) \right\|_2^2}_{\text{regularization term}}\, ,
\]
where $\vb{g}_{\mathrm{obs}}$ are the observed gravity data, $g(\vb{m})$ is the predicted gravity using Equation~\ref{eq:forward}, $\vb{W}_d$ is a diagonal matrix of inverse variance (data uncertainty), and $\lambda$ is the Levenberg--Marquart regularization parameter.  Our inital guess is given by the column vector $\vb{m}_0$.

To determine $\vb{m}$, we have to solve iteratively, with each model step being
\[
  \Delta \vb{m} = (\vb{J}\T \vb{W}_d^2 \vb{J} + \lambda^2 \vb{I})^{-1} 
  \vb{J}\T \vb{W}_d^2 \, (\vb{g}_{\mathrm{obs}} - g(\vb{m}_0))\, .
\]

\paragraph{Role of $\lambda$.} Small $\lambda$ approaches Gauss--Newton (fast, but can be unstable if $J^\top J$ is ill-conditioned). Large $\lambda$ yields smaller, more stable steps (gradient-descent-like), helpful far from the solution or with strong trade-offs.

\begin{questions}
    \question{Jacobian}
    Before we play with the buried sphere program, let's derive derivatives we need for the Jacobian matrix.

    The Jacobian for the buried sphere requires three derivatives:
    \[
        \vb{J} =
        \begin{bmatrix} 
            \displaystyle\pdv{g}{z_c} & \displaystyle\pdv{g}{a} & \displaystyle\pdv{g}{\Delta\rho}
        \end{bmatrix}\, .
    \]
    Derive the derivatives from Equation~\ref{eq:forward}.
    \answerbox[460pt]
\end{questions}

\section*{Exploration}

The questions below relate to the \verb+buried_sphere_gui.py+ program. While the questions are about the demo, the answers are largely independent of both the problem and program.

\clearpage
\begin{questions}[resume]
    \question{Grid Search}
    What is a grid search? Is it a forward or inverse driven solution? What does it tell us about the misfit? How do we determine the model parameters?
    \answerbox

    \question{Inversion}
    How does inversion approach the problem differently from the grid search?  What path does it take to the answer?
    \answerbox

    \question{Comparison}
    How does the result of the inversion compare with that of the grid search? If the two solutions do not match, why?
    \answerbox

    \clearpage
    \question{Non-uniqueness}
    The problem we have posed is non-unique even though we use a single "true" model to generate the data.  Why is the solution non-unique?
    \answerbox

    \question{2D versus 3D}
    Let's switch from 3 parameters to 2.  The 2 parameter result is a better constrained, why?  How is this improvement in constraint reflected in the misfit surfaces?
    \answerbox

    \question{Noise}
    How does noise affect the solution via inversion?  How does noise contribute to non-uniqueness?

    For this demonstration, repeat the inversion with 0\%, 5\%, and 10\% Gaussian noise.
    \answerbox
\end{questions}

\section*{Key Ideas}

\textbf{Grid search} is a brute-force forward modelling approach to systematically sample the model space with misfit is computed at every combination of model parameters. It can identify the global minimum within the explored model domain, but it becomes computationally expensive as the number of model parameters increases and requires a sufficiently dense sampling of the model space to accurately resolve the misfit surface. For many large difficult problems found in geophysics, a grid search is unfeasible.

\textbf{Inversion} uses information about the local shape of the misfit surface, typically through the Jacobian matrix, to move iteratively toward a low-misfit solution. While each iteration is computationally more expensive than a single forward calculation, inversion generally requires far fewer model evaluations than a grid search.  Best of all, inversion remains practical for geophysical problems with many parameters (tens to millions). The trade-off is that inversion explores only a small portion of model space and may converge to a local minimum rather than the global minimum.

\begin{figure}[htb]
\centering

\begin{minipage}{0.48\textwidth}
\centering
\textbf{Grid Search}
\vspace{1em}

\begin{tikzpicture}[
    node distance=0.9cm,
    block/.style={
        rectangle,
        draw,
        align=center,
        minimum width=3.2cm,
        minimum height=0.9cm
    }
]

\node[block] (model) {Choose\\Model Domain};

\node[block, below=of model] (forward) {Forward Model};

\node[block, below=of forward] (misfit) {Compute\\Misfit};

\node[block, below=of misfit] (surface) {Construct\\Misfit Surface};

\node[block, below=of surface] (minimum) {Locate Global\\Minimum};

\draw[->] (model) -- (forward);
\draw[->] (forward) -- (misfit);
\draw[->] (misfit) -- (surface);
\draw[->] (surface) -- (minimum);

\end{tikzpicture}
\end{minipage}
\hfill
\begin{minipage}{0.48\textwidth}
\centering
\textbf{Gradient-Based Inversion}
\vspace{1em}

\begin{tikzpicture}[
    node distance=0.9cm,
    block/.style={
        rectangle,
        draw,
        align=center,
        minimum width=3.2cm,
        minimum height=0.9cm
    },
    decision/.style={
        diamond,
        draw,
        aspect=2,
        align=center,
        inner sep=1pt
    }
]

\node[block] (start) {Choose\\Initial Model};

\node[block, below=of start] (forward) {Forward Model};

\node[block, below=of forward] (jac) {Jacobian};

\node[block, below=of jac] (update) {Compute Model\\Update};

\node[decision, below=1.2cm of update]
      (conv) {Misfit\\Tolerance\\Met?};

\node[block, below=1.2cm of conv]
      (done) {Final Model};

\draw[->] (start) -- (forward);
\draw[->] (forward) -- (jac);
\draw[->] (jac) -- (update);
\draw[->] (update) -- (conv);

\draw[->] (conv) -- node[right]{Yes} (done);

% \draw[->]
%   (conv.west)
%   -- ++(-1.7,0)
%   node[midway,left]{No}
%   |- (forward.west);

% \node[left=1.9cm of conv]
%       {\small Iterate};
\coordinate (loopA) at ($(conv.west)+(-1.7,0)$);

\draw[->] (conv.west) -- (loopA);
\draw[->] (loopA) |- (forward.west);

\node[below] at ($(conv.west)!0.5!(loopA)$) {No};

\node[
    rotate=90,
    anchor=center
] at ($(loopA)!0.5!(forward.west)$)
{\small Iterate};

\end{tikzpicture}
\end{minipage}

\caption{Comparison of grid-search and gradient-based inversion approaches.}
\end{figure}

\end{activity}